\originalchapter{回归与预测}
\originalsection{总体的最佳直线}
在需求与价格、温度与气压等数据中, 我们观察成对变量$(X_i,Y_i)$, 希望用$X_i$说明$Y_i$的变化. 经济学中的需求可近似写成$D_i=a+bp_i+\varepsilon_i$. 理想气体关系$PV=nRT$取对数后成为$\log P=\log(nR)-\log V+\log T$, 因而可以用$\log V,\log T$作为解释变量建立线性模型. ``线性''指对待估系数线性, 不要求对原始物理变量线性.

\begin{definition}[总体线性回归]
设实值变量$X,Y$二阶矩有限, $\Var(X)>0$. $Y$对$X$的总体线性回归是使$\E[(Y-a-bX)^2]$最小的随机变量$a+bX$.
\end{definition}
\begin{proposition}
最优系数唯一, 且
\[
 b=\frac{\Cov(X,Y)}{\Var(X)},\qquad a=\E Y-b\E X.
\]
记$\varepsilon=Y-a-bX$, 则$\E\varepsilon=0$, $\Cov(X,\varepsilon)=0$. 反之, 若某种分解满足这两个正交条件, 则其直线就是总体线性回归.
\end{proposition}
\begin{proof}
先固定$b$, 用$\E[(W-a)^2]=\Var(W)+(\E W-a)^2$知最优$a=\E Y-b\E X$. 代入得
\[
 \E[(Y-a-bX)^2]=\Var(Y)-2b\Cov(X,Y)+b^2\Var(X).
\]
因$\Var(X)>0$, 配方后的唯一最小点如上. 正交条件由代入直接得到. 反之, 对任意增量$u,v$, 有
\[
 \E[(\varepsilon-u-vX)^2]
 =\E\varepsilon^2+\E[(u+vX)^2],
\]
因为交叉项$\E[\varepsilon(u+vX)]$为零. 最后一项非负, 且在$\Var(X)>0$下仅当$u=v=0$取零.
\end{proof}
独立且中心的误差足以保证正交, 条件均值$\E[\varepsilon\mid X]=0$也足够. 但零协方差较弱, 不保证独立, 也不保证条件均值为零. 最佳线性近似可能不同于完整回归函数$\E[Y\mid X]$, 后者可以是非线性曲线.

\begin{remark}[回归现象]
这是为用户考纲补充的总体解释. 若还假设$\Var(Y)>0$, 记$\sigma_X,\sigma_Y$为两总体标准差, $r=\Cov(X,Y)/(\sigma_X\sigma_Y)$为相关系数. 对标准化变量$X^*=(X-\E X)/\sigma_X$, $Y^*=(Y-\E Y)/\sigma_Y$, 前述公式给出最佳线性预测$\widehat{Y^*}=rX^*$. Cauchy--Schwarz不等式给出$|r|\leq1$. 当$|r|<1$时, 一个距离自身均值$|x|$个标准差的解释变量, 所对应的最佳线性预测只距响应均值$|rx|$个标准差. 这就是向均值回归的含义. 它讨论预测位置, 不声称每个个体的下一次观测都更接近均值, 更不证明时间变化的因果机制. 若$(X,Y)$联合正态, 则条件均值恰为这条直线, 条件方差为$\sigma_Y^2(1-r^2)$; 一般分布中这里只保证最佳线性预测, 不能自动把它认作条件均值.
\end{remark}

\originalsection{样本最小二乘与残差}
\begin{definition}[简单回归的最小二乘]
对观测对$(X_i,Y_i)$, 最小二乘估计$(\hat a,\hat b)$最小化$\sum_i(Y_i-a-bX_i)^2$. 拟合值为$\hat Y_i=\hat a+\hat bX_i$, 残差为$e_i=Y_i-\hat Y_i$.
\end{definition}
\begin{proposition}
若$S_{XX}=\sum_i(X_i-\bar X)^2>0$, 则
\[
 \hat b=\frac{\sum_i(X_i-\bar X)(Y_i-\bar Y)}{S_{XX}},\qquad
 \hat a=\bar Y-\hat b\bar X.
\]
并有$\sum_ie_i=0$, $\sum_iX_ie_i=0$.
\end{proposition}
\begin{proof}
平方和对$a,b$求偏导得到$\sum_i(Y_i-a-bX_i)=0$与$\sum_iX_i(Y_i-a-bX_i)=0$. 第一式给出$a=\bar Y-b\bar X$, 代入第二式得到斜率公式. $S_{XX}>0$使二次目标的Hessian正定, 因而驻点是唯一全局最小点. 两个残差条件就是原正规方程.
\end{proof}
课件的连续图示依次显示散点、真实直线、估计直线、两条直线并列以及观测点到拟合直线的竖直距离. 其中真实直线是一种总体关系, 拟合直线是随机样本算出的关系, 二者通常不同. 最小二乘最小化这些竖直残差的平方, 不最小化点到直线的垂直几何距离. 截距保证拟合直线经过$(\bar X,\bar Y)$.

\begin{example}
对三点$(0,1),(1,2),(2,2)$求拟合直线与残差平方和.
\end{example}
\begin{solution}
$\bar X=1$, $\bar Y=5/3$, $S_{XX}=2$, $S_{XY}=1$, 所以$\hat b=1/2$, $\hat a=7/6$. 拟合值为$(7/6,5/3,13/6)$, 残差为$(-1/6,1/3,-1/6)$, 平方和为$1/6$. 残差和为0, 与$X$的内积也为$1/3-2/6=0$, 可直接验证正规方程.
\end{solution}
独立同分布观测对满足所需二阶矩时, 大数定律使样本方差与协方差趋于总体量, 因总体$\Var(X)>0$, 上述比值给出$\hat b\pto b$, $\hat a\pto a$. 此一致性不要求误差正态; 精确有限样本推断则需要后面的更强模型.

\originalsection{多元回归与正交投影}
设第$i$个观测有列向量$x_i\in\R^p$, 通常第一坐标为1以容纳截距. 设计矩阵$X\in\R^{n\times p}$的第$i$行为$x_i^\top$, 响应向量$Y\in\R^n$, 模型为
\[
 Y=X\beta+\varepsilon,\qquad\beta\in\R^p.
\]
注意设计矩阵$X$与前一节的标量解释变量同名, 对象类型已发生变化. 最小二乘选择使$\norm{Y-Xb}^2$最小的$b$.
\begin{theorem}[满列秩正规方程]
若$\operatorname{rank}(X)=p$, 则$X^\top X$正定, 最小二乘解唯一并等于
\[
 \hat\beta=(X^\top X)^{-1}X^\top Y.
\]
帽子矩阵$H=X(X^\top X)^{-1}X^\top$是到$\operatorname{col}(X)$的正交投影. 拟合值$\hat Y=HY$, 残差$e=(I-H)Y$, 满足$X^\top e=0$.
\end{theorem}
\begin{proof}
对非零$b$, $b^\top X^\top Xb=\norm{Xb}^2>0$. 求导给出正规方程$X^\top X\hat\beta=X^\top Y$, 可逆性给出解. $H^\top=H$, $H^2=H$, 且$HX=X$, 故它固定列空间并将其正交补送到0. 还可将目标写为
\[
 \norm{Y-Xb}^2=\norm{(I-H)Y}^2+\norm{HY-Xb}^2,
\]
两项正交, 因而拟合值确为最佳投影.
\end{proof}
如果列之间线性相关, 参数的某些组合无法从$X\beta$分辨, 不能继续写逆矩阵公式. 计算上即使满秩, 严重共线性也会使$X^\top X$条件数很大, 从而放大系数估计误差; 这与数学上的唯一性是不同层次的问题.

\originalsection{正态固定设计模型}
以下假设$X$固定、满列秩, $n>p$, 以及$\varepsilon\sim N_n(0,\sigma^2I_n)$, $\sigma^2>0$. 这同时要求误差中心、同方差、相互独立和正态. 若设计随机, 可以在条件于整个$X$后假设$\varepsilon\mid X\sim N_n(0,\sigma^2I_n)$; 仅有逐个误差与某个解释变量不相关还不够.

\begin{theorem}[最小二乘的精确分布]
在上述条件下,
\[
 \hat\beta\sim N_p(\beta,\sigma^2(X^\top X)^{-1}),\qquad
 \frac{\operatorname{RSS}}{\sigma^2}\sim\chi_{n-p}^2,
 \qquad \operatorname{RSS}=\norm{Y-X\hat\beta}^2,
\]
而且$\hat\beta$与$\operatorname{RSS}$独立. $s^2=\operatorname{RSS}/(n-p)$是$\sigma^2$的无偏估计. 同时
\[
 \E\norm{\hat\beta-\beta}^2=\sigma^2\tr((X^\top X)^{-1}),\quad
 \E\operatorname{RSS}=(n-p)\sigma^2,\quad
 \E\norm{\hat Y-X\beta}^2=p\sigma^2.
\]
\end{theorem}
\begin{proof}
$\hat\beta-\beta=(X^\top X)^{-1}X^\top\varepsilon$, 是正态向量的线性变换, 协方差按矩阵相乘得到. $e=(I-H)\varepsilon$, $I-H$为秩$n-p$的正交投影. 在适合列空间与正交补的标准正交基中, 投影噪声和残差噪声取自独立标准正态坐标, 于是得到卡方分布与独立性. 平方范数的期望等于协方差的迹, 故三条风险公式分别来自$(X^\top X)^{-1}$、$I-H$、$H$的迹.
\end{proof}
正态对数似然关于$\beta$只通过$-\operatorname{RSS}/(2\sigma^2)$变化, 所以$\hat\beta$也是MLE. 若方差未知, 其MLE为$\operatorname{RSS}/n$, 不同于无偏估计$s^2$. 课件把$\E\norm{Y-X\hat\beta}^2$称为预测误差; 它实际是训练残差. 对独立新响应向量$Y^{\rm new}=X\beta+\varepsilon^{\rm new}$, 同一设计上的预测风险为
\[
 \E\norm{Y^{\rm new}-\hat Y}^2=n\sigma^2+p\sigma^2,
\]
因为新噪声与估计误差独立. 区分训练残差与未来预测误差可以看清过拟合.

\originalsection{系数、均值与新观测的区间}
\begin{proposition}[线性对比的$t$枢轴]
对固定非零$c\in\R^p$, 记$v_c=c^\top(X^\top X)^{-1}c>0$. 则
\[
 \frac{c^\top\hat\beta-c^\top\beta}{s\sqrt{v_c}}\sim t_{n-p}.
\]
\end{proposition}
\begin{proof}
分子除以$\sigma\sqrt{v_c}$是标准正态; $s^2/\sigma^2$是独立的$\chi_{n-p}^2/(n-p)$, 因而比值按定义服从Student分布.
\end{proof}
取$c$为第$j$个坐标向量, 得到检验$H_0:\beta_j=0$的统计量$\hat\beta_j/[s\sqrt{((X^\top X)^{-1})_{jj}}]$. 双侧水平$\alpha$用$t_{n-p,1-\alpha/2}$校准. 同一枢轴反演出置信区间
\[
 c^\top\hat\beta\ \pm\ t_{n-p,1-\alpha/2}s\sqrt{v_c}.
\]
取$c=x_0$便得到设计点$x_0$处均值$x_0^\top\beta$的区间. 若要预测独立新观测$Y_0=x_0^\top\beta+\varepsilon_0$, 新噪声的方差还要加上$\sigma^2$, 所以预测区间为
\[
 x_0^\top\hat\beta\ \pm\ t_{n-p,1-\alpha/2}s
 \sqrt{1+x_0^\top(X^\top X)^{-1}x_0}.
\]
其枢轴精确性来自$\varepsilon_0$与原样本独立且正态. 同一点的预测区间比均值区间宽, 因为未来个体自身仍会波动.

\begin{example}
沿用$(0,1),(1,2),(2,2)$, 在$x_0=1$处比较标准误.
\end{example}
\begin{solution}
设计列为截距与$x$, 故$n=3,p=2$, $s^2=(1/6)/(3-2)=1/6$. 中心点均值的$v_{x_0}=1/3+(x_0-\bar X)^2/S_{XX}=1/3$. 均值估计为$5/3$, 标准误$\sqrt{1/18}$; 新观测标准误$\sqrt{(1/6)(1+1/3)}=\sqrt{2/9}$. 两个区间都用$t_1$分位数, 因自由度只有1, 不能拿1.96替代.
\end{solution}

\originalsection{联合显著性、F检验与方差分析}
同时检查$k$个系数时, 若分别使用水平$\alpha$, ``至少一个误拒绝''的概率可能超过$\alpha$. Bonferroni规则把每个检验水平设为$\alpha/k$, 只要一个拒绝便拒绝联合原假设. 联合原假设下由并集界
\[
 \PP\Big(\bigcup_{j\in S}\{\text{第}j\text{项拒绝}\}\Big)
 \leq\sum_{j\in S}\alpha/k=\alpha.
\]
不要求各检验独立. 对一个预先指定的整体线性约束, F检验则直接测量联合偏离.

\begin{definition}[Fisher分布]
若$U\sim\chi_k^2$, $V\sim\chi_q^2$独立, 则$(U/k)/(V/q)$服从$F_{k,q}$.
\end{definition}
\begin{theorem}[一般线性约束检验]
设$G\in\R^{k\times p}$行满秩, $\lambda\in\R^k$固定, $H_0:G\beta=\lambda$. 则
\[
 F=\frac{(G\hat\beta-\lambda)^\top
 [G(X^\top X)^{-1}G^\top]^{-1}(G\hat\beta-\lambda)}{ks^2}
 \sim F_{k,n-p}
\]
在原假设下精确成立, 用其$1-\alpha$分位数作上尾阈值.
\end{theorem}
\begin{proof}
设$C=G(X^\top X)^{-1}G^\top$. 因$G$满行秩且中间矩阵正定, $C$正定. 在原假设下,$C^{-1/2}(G\hat\beta-\lambda)/\sigma$为$k$维标准正态, 故分子除以$\sigma^2$为$\chi_k^2$. 它是$\hat\beta$的函数, 与$s^2$独立. 分母的卡方表示随即给出F分布.
\end{proof}
分子也等于约束模型与完整模型的残差平方和差$\operatorname{RSS}_0-\operatorname{RSS}$, 因而
\[
 F=\frac{(\operatorname{RSS}_0-\operatorname{RSS})/k}{\operatorname{RSS}/(n-p)}.
\]
确实, 令$A=X^\top X$. 目标为$\operatorname{RSS}+(b-\hat\beta)^\top A(b-\hat\beta)$. 用乘子最小化第二项且$Gb=\lambda$, 得约束解
\[
 \hat\beta_0=\hat\beta-A^{-1}G^\top(GA^{-1}G^\top)^{-1}(G\hat\beta-\lambda).
\]
代回二次项便得到分子. 这说明F检验比较为满足约束而增加的误差与剩余噪声水平.

若模型含截距, 置$\operatorname{TSS}=\sum_i(Y_i-\bar Y)^2$, $\operatorname{ESS}=\sum_i(\hat Y_i-\bar Y)^2$. 正交性给出$\operatorname{TSS}=\operatorname{ESS}+\operatorname{RSS}$. 当$p>1$时, 检验所有非截距系数为零可写成
\[
 F=\frac{\operatorname{ESS}/(p-1)}{\operatorname{RSS}/(n-p)}.
\]
这是回归方差分析的平方和分解; 在原假设下两个平方和除以$\sigma^2$独立, 自由度分别为$p-1$与$n-p$. 决定系数$R^2=1-\operatorname{RSS}/\operatorname{TSS}$在$\operatorname{TSS}>0$时定义, 衡量训练样本中的解释比例; 加变量会提高或保持$R^2$, 不保证未来预测改善.

\originalsection{杠杆、残差与模型诊断}
帽子矩阵对角元$h_{ii}=x_i^\top(X^\top X)^{-1}x_i$称为杠杆值. 投影性质给出$0\leq h_{ii}\leq1$, 且$\sum_i h_{ii}=\tr H=p$. 第一条可由$h_{ii}=\norm{He_i^{\rm coord}}^2$及投影不增范数得到. 简单回归含截距时
\[
 h_{ii}=\frac1n+\frac{(X_i-\bar X)^2}{S_{XX}},
\]
所以远离解释变量中心的点杠杆更高. 高杠杆是设计位置的性质, 不等于响应一定异常.

误差$\varepsilon$不可直接观察, 残差$e=(I-H)\varepsilon$虽为正态, 却满足
\[
 \Cov(e)=\sigma^2(I-H),\qquad\Var(e_i)=\sigma^2(1-h_{ii}),\qquad
 \Cov(e_i,e_j)=-\sigma^2h_{ij}\quad(i\ne j).
\]
它们一般不独立也不同方差. 因而对残差作QQ图可帮助诊断, 但不能把残差当独立同分布样本直接套普通KS或Lilliefors表. 合适的模拟诊断应在同一设计下生成模型数据、重新拟合, 并校准整个残差统计量. 回归揭示变量之间的关联; 因果解释需要实验设计或额外识别条件, 不来自显著系数本身.

\originalsection{秩亏时仍可拟合什么}
当$k=\operatorname{rank}(X)<p$, 若$v\in\ker X$非零, 则$X(\beta+v)=X\beta$, 所以分布无法区分$\beta$与$\beta+v$. 系数不再整体可识别, 但均值向量$\mu=X\beta\in\operatorname{col}(X)$仍可识别.
\begin{proposition}[秩亏设计]
令$H$为到$\operatorname{col}(X)$的正交投影, 可写成$H=X(X^\top X)^\dagger X^\top$, 其中$\dagger$为Moore--Penrose广义逆. 所有最小二乘解给出同一拟合值$HY$. 正态同方差模型下
\[
 \norm{(I-H)Y}^2/\sigma^2\sim\chi_{n-k}^2,
 \qquad HY\ \text{与}\ \norm{(I-H)Y}^2\ \text{独立}.
\]
其中$k=n$时$\chi_0^2$约定表示集中于0的点质量. 若$k<n$, $\norm{(I-H)Y}^2/(n-k)$仍是方差无偏估计.
\end{proposition}
\begin{proof}
取列空间的标准正交基组成$U\in\R^{n\times k}$, 则$H=UU^\top$. 对任意$Xb$, 正交分解证明最小残差只能是$(I-H)Y$, 可达到的拟合值为$HY$. 用$X$的奇异值分解核对广义逆表达式. 高斯噪声在列空间与正交补的坐标独立, 残差有$n-k$个标准正态坐标, 得分布与独立性.
\end{proof}
当$k=n$时训练残差恒为0, 没有剩余自由度估计噪声方差. 这不说明噪声不存在. 某些系数组合仍可识别: $c^\top\beta$可识别当且仅当$c$属于$X$的行空间, 因为它必须对所有$v\in\ker X$满足$c^\top v=0$. 类似地, 新点预测$x_0^\top\beta$只有在$x_0$属于行空间时由训练设计识别.

\originalsection{高维变量选择与惩罚}
当$p$大于$n$或接近$n$, 即使许多变量无关, 仍可用大量自由参数逼近训练数据. 课程提出稀疏假设: 只有少数系数非零, 但其位置未知. 对子集$S\subset\{1,\ldots,p\}$, 用子矩阵$X_S$拟合$\hat\beta_S$, 比较
\[
 \operatorname{RSS}_S+\lambda|S|,\qquad\lambda>0.
\]
残差项鼓励拟合, 惩罚项限制模型复杂度. 单比较训练残差会偏向最大模型, 因为扩大拟合空间不会增加残差.

在固定候选设计、已知同方差高斯噪声下, Mallows的$C_p$可写为
\[
 C_p(S)=\operatorname{RSS}_S/\sigma^2-n+2k_S,
 \qquad k_S=\operatorname{rank}(X_S).
\]
常数项与缩放不影响选择, 故等价于最小化$\operatorname{RSS}_S+2\sigma^2k_S$. 对满列秩子模型$k_S=|S|$. 其意义可由直接风险计算说明: 对任意真实均值$\mu$, 子模型投影$H_S$满足
\[
 \E\operatorname{RSS}_S=\norm{(I-H_S)\mu}^2+(n-k_S)\sigma^2,
\]
而均值预测风险为$\E\norm{H_SY-\mu}^2=\norm{(I-H_S)\mu}^2+k_S\sigma^2$. 所以$\operatorname{RSS}_S-n\sigma^2+2k_S\sigma^2$是后者的无偏估计. 这里子模型须预先固定; 从很多模型中挑最小者后, 不能再把选中模型的误差估计直接视为无选择偏差.

已知方差正态似然的AIC为$\operatorname{RSS}_S/\sigma^2+2k_S$加共同常数, BIC为$\operatorname{RSS}_S/\sigma^2+k_S\log n$加共同常数. 所以对应$\lambda=2\sigma^2$与$\sigma^2\log n$. 课件用$\hat\sigma^2$代替作为实践形式; 若每个模型分别最大化未知方差, 精确似然形式则是$n\log(\operatorname{RSS}_S/n)+2(k_S+1)$或$n\log(\operatorname{RSS}_S/n)+(k_S+1)\log n$. 因此``AIC等于残差加固定惩罚''须注明方差处理方式.

\begin{definition}[$\ell_0$与Lasso]
$\norm b_0=\sum_j\ind_{\{b_j\ne0\}}$记非零坐标数, 它不是通常意义下的范数. 最佳子集问题可写成最小化$\norm{Y-Xb}^2+\lambda\norm b_0$. Lasso用$\ell_1$惩罚替换它:
\[
 \hat\beta^{\rm L}\in\argmin_{b\in\R^p}
 \left\{\norm{Y-Xb}^2+\lambda\sum_{j=1}^p|b_j|\right\}.
\]
\end{definition}
最佳子集通常要面对$2^p$个候选子集, 不是课件所写的$2^n$. Lasso目标是凸函数, 且$\lambda>0$使它随$\norm b\to\infty$趋于无穷, 所以存在最小点. 设计秩亏时系数解未必唯一. 截距在实践中常不惩罚, 变量通常先标准化; 两种选择均须在定义计算程序时固定.

\begin{example}
若$X^\top X=I_p$, 求Lasso解.
\end{example}
\begin{solution}
记$z=X^\top Y$, 忽略与$b$无关的常数, 目标为$\sum_j\{(b_j-z_j)^2+\lambda|b_j|\}$. 对$b_j>0$求导得到$b_j=z_j-\lambda/2$, 仅在$z_j>\lambda/2$可行; 对负半轴得到$b_j=z_j+\lambda/2$, 仅在$z_j<-\lambda/2$可行; 其余情形最小点在0. 因而
\[
 \hat b_j=\operatorname{sign}(z_j)(|z_j|-\lambda/2)_+.
\]
随着$\lambda$增大, 系数收缩且依次成为0, 这解释了原课件的系数路径图. 阈值$\lambda/2$取决于本文未在平方损失前乘$1/2$的约定.
\end{solution}
一般设计的路径可以改变斜率, 个别坐标也可能进入后再退出模型; 不能把正交例的单调坐标行为当作一般结论. 惩罚参数常用验证数据或交叉验证选择. 要从小预测误差进一步推断真实支持集恢复, 还需设计条件、信号大小与噪声控制, 并非选择任意$\lambda$就能保证. 课程把这些问题留给高维统计课程.

\originalsection{非参数回归与局部平均}
线性条件均值并非总合理. 设$Y_i=f(X_i)+\varepsilon_i$, $\E[\varepsilon_i\mid X_i]=0$, 希望估计$f(x)=\E[Y_i\mid X_i=x]$. 参数假设$f(x)=a+bx$, $a+bx+cx^2$或$e^{a+bx}$把问题缩成有限参数估计. 前两式对系数线性, 第三式一般是非线性回归; 不能把所有非线性均值函数都直接套线性矩阵公式.

\begin{definition}[局部常数估计]
对一维设计、带宽$h>0$, 令$I_x=\{i:|X_i-x|<h\}$, $N_x=|I_x|$. 定义
\[
 \hat f_{n,h}(x)=
 \begin{cases}N_x^{-1}\sum_{i\in I_x}Y_i,&N_x>0,\\0,&N_x=0.\end{cases}
\]
空窗取0是课程的一种约定, 不是有数据支持的估计; 实际可选择扩大窗口或报告缺少局部观测.
\end{definition}
思路是只在$x$附近以常数$f(x)$近似$f$. 课程示意图取$x=.6,h=.1$, 对$.5<X_i<.7$的响应取平均, 图中得到约$.27$. 区间的宽度是$2h$, 而不是$h$. 如果$h$太小, 可用观测少且曲线随个别点跳动; 如果$h$太大, 局部结构被平均掉.

\begin{proposition}[条件偏差与方差]
条件于设计, 假设各误差独立、条件均值0、条件方差$\sigma^2$, 且$N_x>0$. 则
\[
 \E[\hat f_{n,h}(x)\mid X]=\frac1{N_x}\sum_{i\in I_x}f(X_i),\qquad
 \Var(\hat f_{n,h}(x)\mid X)=\frac{\sigma^2}{N_x}.
\]
若在该窗口内$|f(t)-f(x)|\leq L|t-x|^s$, $0<s\leq1$, 则条件均方误差不超过$L^2h^{2s}+\sigma^2/N_x$.
\end{proposition}
\begin{proof}
条件期望对平均逐项作用, 条件独立使平均的方差为$N_x\sigma^2/N_x^2$. 偏差是$f(X_i)-f(x)$的平均, 绝对值不超过$Lh^s$. 使用均方误差等于偏差平方加方差即可.
\end{proof}
例如课程用$n=100,f(x)=x(1-x)$比较$h=.005,1,.2$. 极窄窗口$.005$使估计受到稀少点的强烈扰动并可能出现空窗; $h=1$在$[0,1]$内几乎把所有样本纳入, 得到接近水平的总体平均, 不能表达抛物线; $.2$在降噪与保留弯曲间折中. 这些图展示同一估计规则的不同带宽, 不证明$.2$总是最优.

若设计在$x$附近有正且连续密度$q(x)$, 局部点数典型规模约为$2nhq(x)$, 方差规模为$1/(nh)$, 偏差平方规模不超过$h^{2s}$. 平衡两者得到$h$的规模$n^{-1/(2s+1)}$, 因此既要$h\to0$以控制偏差, 又要$nh\to\infty$以积累局部观测. 这是含密度与平滑条件的尺度分析; 若光滑性更高, 对称窗口的消去效应或局部多项式可能给出更好偏差阶, 本式没有声称覆盖这些改进.

平滑性未知时可用交叉验证: 逐个删去观测, 用其余数据计算$\hat f_{-i,h}(X_i)$, 比较
\[
 \operatorname{CV}(h)=\frac1n\sum_i\{Y_i-\hat f_{-i,h}(X_i)\}^2.
\]
必须先固定空窗处理规则并在全部候选带宽中保持一致. 删去自身避免极小带宽仅靠$Y_i$预测$Y_i$而得到虚假的小训练误差. 这里选择的是预测表现良好的带宽, 对其统计最优性仍需另行条件与证明.
